\documentclass[10pt]{article}

% Specify the margins and text size.
\setlength{\textwidth}{6.7in}
\setlength{\textheight}{9.5in}
\setlength{\oddsidemargin}{0pt}
\setlength{\evensidemargin}{0pt}
\setlength{\topmargin}{0pt}
\setlength{\hoffset}{-.1in}
\setlength{\voffset}{-1in}

\setlength{\parskip}{5pt}
\setlength{\parindent}{0pt}

% Load some fonts and symbol packages
\usepackage{latexsym}
\usepackage{pifont}       % contains 'star' symbol for counterinsurgency handbook title
\usepackage{yfonts} 
\usepackage{amsmath}
\usepackage{amsfonts}

\usepackage{graphicx}     % actually, this is loaded by pstricks
\usepackage[T1]{fontenc}
\usepackage{ifthen}
\usepackage{pstricks,pst-grad,pst-text,pst-node,multido,pst-plot,calc,pst-3dplot}
%\usepackage[all]{xy}
%\usepackage{animate}

% The hyperref package inserts links into the pdf.
\definecolor{MyLinkColor}{rgb}{.1,.2,1}
\definecolor{MyCiteColor}{rgb}{.1,1,.2}
\definecolor{MyURLColor}{rgb}{.4,.4,.4}
\usepackage[backref=true,pagebackref=false,hyperindex,colorlinks=true,
  linkcolor=MyLinkColor,urlcolor=MyURLColor]{hyperref}


% The tweaklist package is something I found on the web.  It provides a simple interface
% for making changes to spacing used in the itemize and enumerate environments.  Comment
% this out if you don't care to use tweaklist.
\usepackage{tweaklist}
\renewcommand{\itemhook}{\setlength{\parskip}{2pt}\setlength{\parsep}%
{1pt}\setlength{\topsep}{0pt}\setlength{\itemsep}{0pt}}

\newcommand{\U}{\underline{\hspace{5pt}}}

\usepackage{listings}
\newcommand{\Z}{\hphantom{0}}

\begin{document}
\pagestyle{empty}
\href{http://www.su.edu}{\includegraphics[height=1.75cm]{sulogo.eps}}
\vspace{-1.69cm}

\hfill {\small
\begin{tabular}{cl}
& Math 207\\
& Introduction to Statistics\\
\end{tabular}
}
\setlength{\baselineskip}{1.05\baselineskip}
\bigskip

\begin{center}
\textbf{\large  Exam II}
\end{center}
Please show your work and any R commands that you use.
\medskip


1. \textbf{Standard Normal Curve}:  Sketch and compute the specified area under the normal curve. 
Write down the R command that you use. (15 points)

\hspace{10pt} a) $z < -1$
\vspace{.4in}

\hspace{10pt} b) $z > 1.75$
\vspace{.5in}

\hspace{10pt} c) $-1 < z < 1.75$
\vspace{.5in}

\hspace{10pt} d) What $z$ is greater than 51\% of all $z$ scores?
\vspace{.5in}

2. \textbf{Normal Curve}:  Suppose that for a particular pet population
subjects' weights average out to 52 pounds with an sd of 5 pounds and have
a normal distribution. (15 points)

\hspace{10pt} a) What percent of the weights were below 47 pounds?
\vspace{.9in}

\hspace{10pt} b) What percent were above 62 pounds?
\vspace{.9in}

\hspace{10pt} c) What percent were between 47 and 62 pounds?
\vspace{.9in}

\hspace{10pt} d) What weight was greater than 51\% of all the weights?
\vfill
\eject
{\ }


%\#7 page 144 (word problem about correlation)
%\#8 page 155 (word problem about correlation)
3. \textbf{Correlation}: (10 points) 

\hspace{20pt} a) For a representative sample of cars, would the correlation between the
age of the car and its gasoline economy (miles per gallon) be positive or negative?  Explain
\vspace{1.5in}

\hspace{20pt} b ) The correlation between gasoline economy and income of owner turns out to 
be positive.  How do you account for this association?
\vspace{1.5in}

4. \textbf{RMS Error for a Line}: (15 points) Consider the following data.
\begin{tabular}{|c|c|}\hline
$x$ & $y$\\\hline
0   & 2\\
1   & 4\\
2   & 3\\\hline
\end{tabular}

\hspace{10pt} a) Make a scatter plot of the data.  
\vspace{2in}

\hspace{10pt} b) An investigator uses the line $y=2\,x + 1$ to predict $y$ from $x$.
Add the line to your sketch.
\medskip

\hspace{10pt} c) Calculate the RMS error for the line:
$\displaystyle\sqrt{\frac{1}{n}\sum\left(\mbox{predicted $y$} - \mbox{actual $y$}\right)^2}$.

\vfill
\eject
{\ }

5. \textbf{Regression}: 
(25 points) For people age 25 and over in the U.S. in 2005, the relationship between
age and educational level (years of schooling completed) can be summarized as
follows. \vspace{-5pt}
\begin{align*}
\mbox{average age}         & \approx \mbox{50 years,}    & \mbox{sd} & \approx \mbox{16 years}\\[3pt]
\mbox{average educational level}  & \approx \mbox{13.2 years,}  & \mbox{sd} & \approx \mbox{3.0 years,
   \hspace{15pt}} r\approx -0.2
\end{align*}\vspace{-5pt}

\hspace{20pt} a) Write down the equation for the regression line  to
predict educational level from age.
\vspace{1.4in}

\hspace{20pt} b) Predict the educational level of a 75 year old subject.
\vspace{1.4in}

\hspace{20pt} c) About 68\% of the 75 year old subjects had educational levels in what range?
\vspace{1.4in}

\hspace{20pt} d) True or false:  For people age 25 and older in the U.S. in 2005, older people tended 
to be less\vspace{-5pt}

\hspace{20pt} \hphantom{a) }  educated than younger ones.  Explain.
\vspace{1.4in}


\hspace{20pt} e) True or false, and explain:  as you get older, you become less educated.  
If this statement is false,\vspace{-4pt}

\hspace{20pt} \hphantom{a) }
what could account for the negative correlation?  
\vfill
\eject
{\ }
    
6. \textbf{$z$-Scores}: (10 points) Test scores on a certain exam average out to 75 with a standard 
deviation of 3.  One student scored 84.   Was that a very high score?  Moderately high?  
Or about average?  Justify your answer.
\vspace{3in}


7. \textbf{NHANESIII Data Set}: (10 points)

\hspace{10pt} a) In the exam data set, what information does the variable \texttt{PEP13E2A} give?
What kind of variable is it?
\vspace{2in}


\hspace{10pt} b) In the lab data set, what information does the variable \texttt{HXPAXTMR} give?
What kind of variable is it?
\vfill
\begin{center}
{\footnotesize\begin{pspicture}(-1.5,0)(7.5,2.2)
\psset{linewidth=0.02, nodesep=2pt, xunit=3}
\rput(1,2){\rnode{Variables}{Variables}}
\rput(0,1){\rnode{Qualitative}{Qualitative}}
\rput(2,1){\rnode{Quantitative}{Quantitative}}
\rput(1.5,0){\rnode{Discrete}{Discrete}}
\rput(2.5,0){\rnode{Continuous}{Continuous}}
\rput(-0.5,0){\rnode{Ordered}{Ordered}}
\rput(0.5,0){\rnode{Unordered}{Unordered}}
\ncline{Variables}{Qualitative}
\ncline{Variables}{Quantitative}
\ncline{Quantitative}{Discrete}
\ncline{Quantitative}{Continuous}
\ncline{Qualitative}{Ordered}
\ncline{Qualitative}{Unordered}
\end{pspicture}}
\end{center}


\eject

%%%%%%%%%%%%%%%%%%%%%%%%%%%%%%%%%%%%% VERSION 2
\href{http://www.su.edu}{\includegraphics[height=1.75cm]{sulogo.eps}}
\vspace{-1.69cm}

\hfill {\small
\begin{tabular}{cl}
& Math 207\\
& Introduction to Statistics\\
\end{tabular}
}
\setlength{\baselineskip}{1.05\baselineskip}
\bigskip

\begin{center}
\textbf{\large  Exam II}
\end{center}
Please show your work and any R commands that you use.
\medskip

1. \textbf{$z$-Scores}: (10 points) Test scores on a certain exam average out to 78 with a standard 
deviation of 2.  One student scored 88.   Was that a very high score?  Moderately high?  
Or about average?  Justify your answer.
\vspace{3in}


2. \textbf{NHANESIII Data Set}: (10 points)

\hspace{10pt} a) In the exam data set, what information does the variable \texttt{PEP8A} give?
What kind of variable is it?
\vspace{1.5in}


\hspace{10pt} b) In the lab data set, what information does the variable \texttt{DMPCREGN} give?
What kind of variable is it?
\vfill
\begin{center}
{\footnotesize\begin{pspicture}(-1.5,0)(7.5,2.2)
\psset{linewidth=0.02, nodesep=2pt, xunit=3}
\rput(1,2){\rnode{Variables}{Variables}}
\rput(0,1){\rnode{Qualitative}{Qualitative}}
\rput(2,1){\rnode{Quantitative}{Quantitative}}
\rput(1.5,0){\rnode{Discrete}{Discrete}}
\rput(2.5,0){\rnode{Continuous}{Continuous}}
\rput(-0.5,0){\rnode{Ordered}{Ordered}}
\rput(0.5,0){\rnode{Unordered}{Unordered}}
\ncline{Variables}{Qualitative}
\ncline{Variables}{Quantitative}
\ncline{Quantitative}{Discrete}
\ncline{Quantitative}{Continuous}
\ncline{Qualitative}{Ordered}
\ncline{Qualitative}{Unordered}
\end{pspicture}}
\end{center}

\eject
{\ }

3. \textbf{Standard Normal Curve}:  Sketch and compute the specified area under the normal curve. 
Write down the R command that you use. (15 points)

\hspace{10pt} a) $z > -1$
\vspace{.6in}

\hspace{10pt} a) $z < -1.2$
\vspace{.6in}

\hspace{10pt} c) $-1.2 < z < -1$
\vspace{.6in}

\hspace{10pt} d) What $z$ is \underline{less than} 30\% of all $z$ scores?
\vspace{1in}

4. \textbf{Normal Curve}:  Suppose that for a particular pet population
subjects' weights average out to 8 pounds with an sd of 2 pounds and have
a normal distribution. (15 points)

\hspace{10pt} a) What percent of the weights were below 12 pounds?
\vspace{1in}

\hspace{10pt} b) What percent were above 10 pounds?
\vspace{1in}

\hspace{10pt} c) What percent were between 10 and 12 pounds?
\vspace{1in}

\hspace{10pt} d) What weight was \underline{less than} 30\% of all the weights?
\vfill
\eject
{\ }

5. \textbf{RMS Error for a Line}: (15 points) Consider the following data.
\begin{tabular}{|c|c|}\hline
$x$ & $y$\\\hline
0   & 3\\
1   & 4\\
2   & 2\\\hline
\end{tabular}

\hspace{10pt} a) Make a scatter plot of the data.  
\vspace{2in}

\hspace{10pt} b) An investigator uses the line $y=5 - x$ to predict $y$ from $x$.
Add the line to your sketch.
\medskip

\hspace{10pt} c) Calculate the RMS error for the line:
$\displaystyle\sqrt{\frac{1}{n}\sum\left(\mbox{predicted $y$} - \mbox{actual $y$}\right)^2}$.
\vspace{2.5in}


6. \textbf{Correlation}: (10 points) 

\hspace{20pt} a) For a representative sample of cars, would the correlation between the
age of the car and its gasoline economy (miles per gallon) be positive or negative?  Explain
\vspace{1.5in}

\hspace{20pt} b ) The correlation between gasoline economy and income of owner turns out to 
be positive?  How do you account for this association?


\vfill
\eject
{\ }

7. \textbf{Regression}: 
(25 points) For people age 25 and over in the U.S. in 2005, the relationship between
age and educational level (years of schooling completed) can be summarized as
follows. \vspace{-5pt}
\begin{align*}
\mbox{average age}         & \approx \mbox{52 years,}    & \mbox{sd} & \approx \mbox{10 years}\\[3pt]
\mbox{average educational level}  & \approx \mbox{12.5 years,}  & \mbox{sd} & \approx \mbox{2.5 years,
   \hspace{15pt}} r\approx -0.3
\end{align*}\vspace{-5pt}

\hspace{20pt} a) Write down the equation for the regression line  to
predict educational level from age.
\vspace{1.4in}

\hspace{20pt} b) Predict the educational level of a 60 year old subject.
\vspace{1.4in}

\hspace{20pt} c) About 95\% of the 60 year old subjects had educational levels in what range?
\vspace{1.4in}

\hspace{20pt} d) True or false:  For people age 25 and older in the U.S. in 2005, older people tended 
to be less\vspace{-5pt}

\hspace{20pt} \hphantom{a) }  educated than younger ones.  Explain.
\vspace{1.4in}


\hspace{20pt} e) True or false, and explain:  as you get older, you become less educated.  
If this statement is false,\vspace{-4pt}

\hspace{20pt} \hphantom{a) }
what could account for the negative correlation?  
\vfill
\eject

%%%%%%%%%%%%%%%%%%%%% VERSION 3
\href{http://www.su.edu}{\includegraphics[height=1.75cm]{sulogo.eps}}
\vspace{-1.69cm}

\hfill {\small
\begin{tabular}{cl}
& Math 207\\
& Introduction to Statistics\\
\end{tabular}
}
\setlength{\baselineskip}{1.05\baselineskip}
\bigskip

\begin{center}
\textbf{\large  Exam II}
\end{center}
Please show your work and any R commands that you use.
\medskip


1. \textbf{RMS Error for a Line}: (15 points) Consider the following data.
\begin{tabular}{|c|c|}\hline
$x$ & $y$\\\hline
0   & 1\\
1   & 3\\
2   & 6\\\hline
\end{tabular}

\hspace{10pt} a) Make a scatter plot of the data.  
\vspace{1.7in}

\hspace{10pt} b) An investigator uses the line $y=1 + x$ to predict $y$ from $x$.
Add the line to your sketch.
\medskip

\hspace{10pt} c) Calculate the RMS error for the line:
$\displaystyle\sqrt{\frac{1}{n}\sum\left(\mbox{predicted $y$} - \mbox{actual $y$}\right)^2}$.
\vspace{1.9in}


2. \textbf{Correlation}: (10 points) 

\hspace{20pt} a) For a representative sample of computers, would the correlation between the
amount of memory and the age (in years) of the computer be positive or negative?  Explain.
\vspace{1.1in}

\hspace{20pt} b ) The correlation between amount of computer memory and the income of the owner turns out to 
be positive.  How do you account for this association?
\vfill
\eject
{\ }

3. \textbf{$z$-Scores}: (10 points) Test scores on a certain exam average out to 70 with a standard 
deviation of 10.  One student scored 85.   Was that a very high score?  Moderately high?  
Or about average?  Justify your answer.
\vspace{2.8in}


4. \textbf{NHANESIII Data Set}: (10 points)

\hspace{10pt} a) In the exam data set, what information does the variable \texttt{PEP5R} give?
What kind of variable is it?
\vspace{2in}


\hspace{10pt} b) In the lab data set, what information does the variable \texttt{PHPFAST} give?
What kind of variable is it?
\vfill
\begin{center}
{\footnotesize\begin{pspicture}(-1.5,0)(7.5,2.2)
\psset{linewidth=0.02, nodesep=2pt, xunit=3}
\rput(1,2){\rnode{Variables}{Variables}}
\rput(0,1){\rnode{Qualitative}{Qualitative}}
\rput(2,1){\rnode{Quantitative}{Quantitative}}
\rput(1.5,0){\rnode{Discrete}{Discrete}}
\rput(2.5,0){\rnode{Continuous}{Continuous}}
\rput(-0.5,0){\rnode{Ordered}{Ordered}}
\rput(0.5,0){\rnode{Unordered}{Unordered}}
\ncline{Variables}{Qualitative}
\ncline{Variables}{Quantitative}
\ncline{Quantitative}{Discrete}
\ncline{Quantitative}{Continuous}
\ncline{Qualitative}{Ordered}
\ncline{Qualitative}{Unordered}
\end{pspicture}}
\end{center}

\eject
{\ }

5. \textbf{Standard Normal Curve}:  Sketch and compute the specified area under the normal curve. 
Write down the R command that you use. (15 points)

\hspace{10pt} a) $z > 0.8$
\vspace{.6in}

\hspace{10pt} a) $z < 2$
\vspace{.6in}

\hspace{10pt} c) $0.8 < z < 2$
\vspace{.6in}

\hspace{10pt} d) What $z$ is greater than 40\% of all $z$ scores?
\vspace{1in}

6. \textbf{Normal Curve}:  Suppose that for a particular pet population
subjects' ages average out to 8 years with an sd of 2 years and have
a normal distribution. (15 points)

\hspace{10pt} a) What percent of the pets were older than 3 years?
\vspace{1in}

\hspace{10pt} b) What percent were younger that 12 years?
\vspace{1in}

\hspace{10pt} c) What percent were between 3 and 12 years?
\vspace{1in}

\hspace{10pt} d) What age was greater than 40\% of all the ages?
\vfill
\eject
{\ }


7. \textbf{Regression}: 
(25 points) For people age 25 and over in the U.S. in 2005, the relationship between
age and educational level (years of schooling completed) can be summarized as
follows. \vspace{-5pt}
\begin{align*}
\mbox{average age}         & \approx \mbox{60 years,}    & \mbox{sd} & \approx \mbox{12 years}\\[3pt]
\mbox{average educational level}  & \approx \mbox{11.5 years,}  & \mbox{sd} & \approx \mbox{3.1 years,
   \hspace{15pt}} r\approx -0.3
\end{align*}\vspace{-10pt}

\hspace{20pt} a) Write down the equation for the regression line  to
predict educational level from age.
\vspace{1.4in}

\hspace{20pt} b) Predict the educational level of a 55 year old subject.
\vspace{1.4in}

\hspace{20pt} c) About 95\% of the 55 year old subjects had educational levels in what range?
\vspace{1.4in}

\hspace{20pt} d) True or false:  For people age 25 and older in the U.S. in 2005, older people tended 
to be less\vspace{-5pt}

\hspace{20pt} \hphantom{a) }  educated than younger ones.  Explain.
\vspace{1.4in}


\hspace{20pt} e) True or false, and explain:  as you get older, you become less educated.  
If this statement is false,\vspace{-4pt}

\hspace{20pt} \hphantom{a) }
what could account for the negative correlation?  
\vfill
\eject


\end{document}

5. \textbf{RMS Error for a Line}: (15 points) Consider the following data.
\begin{tabular}{|c|c|}\hline
$x$ & $y$\\\hline
0   & 3\\
1   & 4\\
2   & 2\\\hline
\end{tabular}

\hspace{10pt} a) Make a scatter plot of the data.  
\vspace{2in}

\hspace{10pt} b) An investigator uses the line $y=5 - x$ to predict $y$ from $x$.
Add the line to your sketch.
\medskip

\hspace{10pt} c) Calculate the RMS error for the line:
$\displaystyle\sqrt{\frac{1}{n}\sum\left(\mbox{predicted $y$} - \mbox{actual $y$}\right)^2}$.
\vspace{2.5in}


6. \textbf{Correlation}: (10 points) 

\hspace{20pt} a) For a representative sample of cars, would the correlation between the
weight of a  car and its gasoline economy (miles per gallon) be positive or negative?  Explain
\vspace{1.5in}

\hspace{20pt} b ) The correlation between gasoline economy and car length turns out to 
be negative.  How do you account for this association?


\vfill
\eject
{\ }



6. (15 points) Suppose that the regression line of $y$ on $x$ 
is \[y=\left(\frac{\sqrt{3}}{2}\right)\,\left(\frac{8}{\sqrt{3}}\right)\,(x-1)\] 
where $\mbox{mean}_x=1$, $\mbox{mean}_y=0$, $\mbox{SD}_x=\sqrt{3}$, $\mbox{SD}_y=8$
and $r=\sqrt{3}/2$.\vspace{3pt}

\hspace{20pt} a) Calculate the RMS error for regression 
$\mbox{RMS}=\mbox{SD}_y\,\sqrt{1-r^2}$.  
\vspace{2in}

\hspace{20pt} b) If $x=4$, then we expect 68\% of the $y$-values to be
between \ \underline{\hspace{60pt}} \ \ and \  \underline{\hspace{60pt}} \ .
\vspace{2in}



\vfill
\eject

(15 points)
